% ==============================================================================
% core_step2_motor.tex — SCI/SICE tutorial core deck, Step 2: motors and mixer
% (exercises 3 and 4)
% SCI/SICE チュートリアル要点版 Step 2: モータとミキサ（実習 3, 4）
% ==============================================================================

% ------------------------------------------------------------------
\begin{frame}{モータの配置と Duty}
  \begin{columns}[c]
    \begin{column}{0.44\textwidth}
      \centering
      \adjustbox{max width=\columnwidth, max height=0.75\textheight}{\documentclass[border=10pt]{standalone}
\usepackage{tikz}
\usetikzlibrary{arrows.meta, positioning, calc}

\begin{document}

\definecolor{cwcolor}{RGB}{0, 120, 200}    % Blue for CW
\definecolor{ccwcolor}{RGB}{220, 60, 20}   % Red for CCW
\definecolor{bodycolor}{RGB}{80, 80, 80}   % Dark gray
\begin{tikzpicture}[
    motor/.style={circle, draw, line width=1.5pt, minimum size=18mm, font=\large\bfseries},
    cw/.style={motor, draw=cwcolor, text=cwcolor},
    ccw/.style={motor, draw=ccwcolor, text=ccwcolor},
    label/.style={font=\small, text=bodycolor},
    arrow/.style={-{Stealth[length=3mm]}, line width=1.2pt},
  ]

  % Body (X-frame)
  \draw[bodycolor, line width=2pt] (-2.5, 2.5) -- (2.5, -2.5);
  \draw[bodycolor, line width=2pt] (2.5, 2.5) -- (-2.5, -2.5);

  % Center body
  \fill[bodycolor, rounded corners=3pt] (-0.6, -0.6) rectangle (0.6, 0.6);
  \node[font=\footnotesize, text=white] at (0, 0) {ESP32};

  % Front direction arrow
  \draw[arrow, bodycolor] (0, 1.0) -- (0, 2.0);
  \node[label, above] at (0, 2.0) {Front};

  % M1 - Front Right (CCW)
  \node[ccw] (m1) at (2.5, 2.5) {M1};
  \node[label] at ($(m1.center)+(45:1.75cm)$) {FR};
  \draw[arrow, ccwcolor] ([shift={(-45:1.3)}]m1.center)
    arc[start angle=-45, end angle=225, radius=1.3cm];
  \node[font=\scriptsize, ccwcolor] at (4.2, 2.5) {CCW};

  % M2 - Rear Right (CW)
  \node[cw] (m2) at (2.5, -2.5) {M2};
  \node[label] at ($(m2.center)+(-45:1.75cm)$) {RR};
  \draw[arrow, cwcolor] ([shift={(45:1.3)}]m2.center)
    arc[start angle=45, end angle=-225, radius=1.3cm];
  \node[font=\scriptsize, cwcolor] at (4.2, -2.5) {CW};

  % M3 - Rear Left (CCW)
  \node[ccw] (m3) at (-2.5, -2.5) {M3};
  \node[label] at ($(m3.center)+(225:1.75cm)$) {RL};
  \draw[arrow, ccwcolor] ([shift={(225:1.3)}]m3.center)
    arc[start angle=225, end angle=495, radius=1.3cm];
  \node[font=\scriptsize, ccwcolor] at (-4.2, -2.5) {CCW};

  % M4 - Front Left (CW)
  \node[cw] (m4) at (-2.5, 2.5) {M4};
  \node[label] at ($(m4.center)+(135:1.75cm)$) {FL};
  \draw[arrow, cwcolor] ([shift={(135:1.3)}]m4.center)
    arc[start angle=135, end angle=-135, radius=1.3cm];
  \node[font=\scriptsize, cwcolor] at (-4.2, 2.5) {CW};

  % Legend
  \node[anchor=west, font=\small] at (-4.5, -4.5) {%
    \textcolor{cwcolor}{\rule{8pt}{8pt}} CW (clockwise) \quad
    \textcolor{ccwcolor}{\rule{8pt}{8pt}} CCW (counter-clockwise)};

\end{tikzpicture}
\end{document}
}
    \end{column}
    \begin{column}{0.53\textwidth}
      \footnotesize
      \begin{itemize}\setlength{\itemsep}{3pt}
        \item モータ ID: 1 = FR（右前）、2 = RR（右後）、3 = RL（左後）、4 = FL（左前）
        \item 対角のモータが同じ向きに回り、反作用トルクを打ち消す（M1・M3 = CCW、M2・M4 = CW）
        \item モータは PWM（パルス幅変調: 電圧を入り切りする時間の割合で平均電圧を変える方式）の Duty 比で駆動する。Duty 0 = 停止、1 = 最大
        \item API: \api{ws::motor\_set\_duty(id, duty)}（id = 1--4、duty = 0.0--1.0）
      \end{itemize}
    \end{column}
  \end{columns}
\end{frame}

% ------------------------------------------------------------------
\begin{frame}[fragile]{実習 3: Duty を決め打ちして 1 つのモータを回す}
  \exgoal{Duty とモータ ID の対応を、回る場所で確かめる}
  \begin{columns}[T]
    \begin{column}{0.47\textwidth}
      {\footnotesize\textbf{\textcolor{sfblue}{手順（机上）}}}\par\vspace{-2pt}
      \footnotesize
      \begin{enumerate}\setlength{\itemsep}{0pt}
        \item \code{sf lesson switch sci2026:3}
        \item \code{sf lesson edit} で右のコードにして保存
        \item \code{sf lesson build}、\code{sf lesson flash}
        \item USB を外し、水平な机に置いて、機体ボタンを 1 回クリックして ARM
      \end{enumerate}
    \end{column}
    \begin{column}{0.51\textwidth}
      \begin{sflisting}[basicstyle=\ttfamily\scriptsize]{user\_code.cpp}
void loop_400Hz(float dt) {
    ws::motor_set_duty(1, 0.10f);
    ws::motor_set_duty(2, 0.00f);
    ws::motor_set_duty(3, 0.00f);
    ws::motor_set_duty(4, 0.00f);
}
      \end{sflisting}
      {\scriptsize \warn{Duty は 0.15 以下}（ファームは制限しない）。\api{ws::arm()} は書かず、ボタンで ARM。ALT/POS は OFF（前々ページの安全ルール）\par}
    \end{column}
  \end{columns}
  \exresult{FR（右前）のモータだけが回る。止めるときはボタンを再クリック（DISARM）}
  \exhelp{3}
\end{frame}

% ------------------------------------------------------------------
\begin{frame}{ミキサ: 推力とトルク指令を 4 つの Duty に配る}
  \begin{columns}[T]
    \begin{column}{0.52\textwidth}
      \centering
      \adjustbox{max width=\columnwidth, max height=0.64\textheight}{\documentclass[border=10pt]{standalone}
\usepackage{tikz}
\usepackage{amsmath}
\usetikzlibrary{arrows.meta, positioning}

\begin{document}

\definecolor{thrustcolor}{RGB}{0, 120, 200}   % Blue - Thrust
\definecolor{rollcolor}{RGB}{220, 50, 30}     % Red - Roll
\definecolor{pitchcolor}{RGB}{40, 160, 40}    % Green - Pitch
\definecolor{yawcolor}{RGB}{120, 80, 150}     % Purple - Yaw
\definecolor{motorcolor}{RGB}{80, 80, 80}     % Gray - Motors
\definecolor{pluscolor}{RGB}{40, 160, 40}      % Green - positive
\definecolor{minuscolor}{RGB}{220, 50, 30}     % Red - negative
\begin{tikzpicture}[
    cmd/.style={
      rectangle, draw=#1, line width=1.5pt,
      fill=#1!10, rounded corners=3pt,
      minimum width=20mm, minimum height=12mm,
      font=\bfseries, text=#1
    },
    mot/.style={
      rectangle, draw=motorcolor, line width=1.5pt,
      fill=motorcolor!8, rounded corners=3pt,
      minimum width=20mm, minimum height=12mm,
      font=\bfseries, text=motorcolor
    },
    plus/.style={-{Stealth[length=2.5mm]}, line width=1.2pt, pluscolor},
    minus/.style={-{Stealth[length=2.5mm]}, line width=1.2pt, minuscolor, dashed},
  ]

  % Command inputs (left side)
  \node[cmd=thrustcolor] (T) at (0, 3)  {T};
  \node[cmd=rollcolor]   (R) at (0, 1)  {Roll};
  \node[cmd=pitchcolor]  (P) at (0, -1) {Pitch};
  \node[cmd=yawcolor]    (Y) at (0, -3) {Yaw};

  % Motor outputs (right side)
  \node[mot] (M1) at (8, 3)  {M1 (FR)};
  \node[mot] (M2) at (8, 1)  {M2 (RR)};
  \node[mot] (M3) at (8, -1) {M3 (RL)};
  \node[mot] (M4) at (8, -3) {M4 (FL)};

  % --- Thrust: all positive ---
  \draw[plus] (T.east) -- (M1.west);
  \draw[plus] (T.east) -- (M2.west);
  \draw[plus] (T.east) -- (M3.west);
  \draw[plus] (T.east) -- (M4.west);

  % --- Roll: -M1,-M2, +M3,+M4 ---
  \draw[minus] (R.east) -- (M1.west);
  \draw[minus] (R.east) -- (M2.west);
  \draw[plus]  (R.east) -- (M3.west);
  \draw[plus]  (R.east) -- (M4.west);

  % --- Pitch: +M1,-M2,-M3,+M4 ---
  \draw[plus]  (P.east) -- (M1.west);
  \draw[minus] (P.east) -- (M2.west);
  \draw[minus] (P.east) -- (M3.west);
  \draw[plus]  (P.east) -- (M4.west);

  % --- Yaw: +M1,-M2,+M3,-M4 ---
  \draw[plus]  (Y.east) -- (M1.west);
  \draw[minus] (Y.east) -- (M2.west);
  \draw[plus]  (Y.east) -- (M3.west);
  \draw[minus] (Y.east) -- (M4.west);

  % Matrix representation (below) — matches ws::motor_mixer() exactly
  % (firmware/workshop/main/ws_internal.hpp resolve_motor_request()):
  % out[i] = T + 0.25*(...)/3.7, a shared-scale approximation, NOT the
  % physical tau/L, kappa mixer used by the real flight controller.
  \node[font=\normalsize, motorcolor, anchor=north] at (4, -5) {%
    $\begin{aligned}
      M_1 &= T + \tfrac{0.25}{3.7}(-R+P+Y) \\[2pt]
      M_2 &= T + \tfrac{0.25}{3.7}(-R-P-Y) \\[2pt]
      M_3 &= T + \tfrac{0.25}{3.7}(+R-P+Y) \\[2pt]
      M_4 &= T + \tfrac{0.25}{3.7}(+R+P-Y)
    \end{aligned}$
  };

  % Legend (right of matrix)
  \node[anchor=west, font=\small, pluscolor] at (8.5, -5.5) {%
    \rule{10pt}{2pt} $+$ positive};
  \node[anchor=west, font=\small, minuscolor] at (8.5, -6.3) {%
    - - - $-$ negative};

\end{tikzpicture}
\end{document}
}
    \end{column}
    \begin{column}{0.45\textwidth}
      \footnotesize
      \begin{itemize}\setlength{\itemsep}{3pt}
        \item 入力: 推力 $T\in[0,1]$、ロール・ピッチ・ヨーのトルク指令 $R,P,Y\in[-1,1]$
        \item 出力: 4 モータの Duty。各式の符号はモータ配置に対応する（例: ピッチ$+$ = 機首上げのとき、前の M1・M4 を強め、後ろの M2・M3 を弱める）
        \item API: \api{ws::motor\_mixer(T, R, P, Y)}
      \end{itemize}
      \vspace{0.1em}
      \[
        d_{\mathrm{FR}} = T + \frac{0.25}{3.7}(-R + P + Y)
      \]
      {\scriptsize 他のモータは符号の組だけが違う。係数 $0.25/3.7$ は、旧世代ファームのミキサから引き継いだ調整定数（物理量ではなく経験的なスケール）\par}
    \end{column}
  \end{columns}
\end{frame}

% ------------------------------------------------------------------
\begin{frame}[fragile]{実習 4: スティックでモータを回す}
  \exgoal{スティック指令がミキサを通って、モータの回転差になることを確かめる}
  \begin{columns}[T]
    \begin{column}{0.47\textwidth}
      {\footnotesize\textbf{\textcolor{sfblue}{手順（机上）}}}\par\vspace{-2pt}
      \footnotesize
      \begin{enumerate}\setlength{\itemsep}{0pt}
        \item \code{sf lesson switch sci2026:4}
        \item \code{sf lesson edit} で右のコードにして保存
        \item \code{sf lesson build}、\code{sf lesson flash}
        \item USB を外して机に置き、ARM。コントローラで操作する
      \end{enumerate}
    \end{column}
    \begin{column}{0.51\textwidth}
      \begin{sflisting}[basicstyle=\ttfamily\scriptsize]{user\_code.cpp}
void loop_400Hz(float dt) {
    float t = ws::rc_throttle();
    float r = ws::rc_roll();
    float p = ws::rc_pitch();
    float y = ws::rc_yaw();
    ws::motor_mixer(t, r, p, y);
}
      \end{sflisting}
      {\scriptsize \warn{スロットルは Duty に直結}。15\% 以下に保つ。ALT/POS は OFF\par}
    \end{column}
  \end{columns}
  \exresult{スティックを倒した向きに応じて、モータの回転に差が出る}
  \exhelp{4}
\end{frame}
